% GENERATED, edit the French unit or the translation table
\mxnewpage\section*{Solutions EV4-T1}\label{EV4-T1-CORR}
\textbf{MCQ:} 1 B, 2 B, 3 B, 4 C, 5 C, 6 B, 7 C, 8 A, 9 A, 10 A, 11 A, 12 D. Each correct answer: 0.5 point; total 6.
\paragraph{Checks of the fluency questions.} Order of operations; sum of fractions; subtraction of exponents; multiplying coefficient 1.2; division by 0.8; product 1.25 times 0.8; solving an affine equation; vertical increment divided by horizontal; mean; angle divided by 360; complement to 1; substitution in brackets.
\subsection*{Exercise 1 -- Choosing a distribution}
\begin{enumerate}
\item $E(X)=2$; $E(Y)=-\dfrac{2}{3}+\dfrac{8}{3}=2$. \textit{Marking: 1 point. Result and basic justification expected; half a point possible for a correct but unfinished approach.}
\item $E(Y^2)=\dfrac{2}{3}+\dfrac{64}{3}=22$, variance 18, against 0 for X. \textit{Marking: 2 points. Method or justification: 1; calculation and conclusion: 1.}
\item Linearity: expectation 4. X is also constant. \textit{Marking: 2 points. Method or justification: 1; calculation and conclusion: 1.}
\end{enumerate}
\subsection*{Exercise 2 -- Discrete threshold and continuous model}
\begin{enumerate}
\item $Q\prime=-64e^{-t}<0$. \textit{Marking: 1 point. Result and basic justification expected; half a point possible for a correct but unfinished approach.}
\item $Q(2)>8$, $Q(3)<8$; decreasing: first index 3. \textit{Marking: 2 points. Method or justification: 1; calculation and conclusion: 1.}
\item Geometric, first term 64, common ratio $e^{-1}$; sum $\dfrac{64(1-e^{-4})}{1-e^{-1}}$. \textit{Marking: 2 points. Method or justification: 1; calculation and conclusion: 1.}
\end{enumerate}
\subsection*{Exercise 3 -- Equation and circle}
\begin{enumerate}
\item $(x-1)^2+(y-2)^2=4$; centre (1;2), radius 2. \textit{Marking: 1 point. Result and basic justification expected; half a point possible for a correct but unfinished approach.}
\item (1;0) and (1;4). \textit{Marking: 1 point. Result and basic justification expected; half a point possible for a correct but unfinished approach.}
\item Vertical radius at the point (1;4), horizontal tangent $y=4$. \textit{Marking: 2 points. Method or justification: 1; calculation and conclusion: 1.}
\end{enumerate}
\paragraph{Calibration.} Target duration: 20--25 min for the fluency questions, 80--85 min for the exercises, 10--15 min for proofreading. These durations are editorial estimates to be measured in class. Accept any correct method; do not penalise twice a numerical error that is propagated consistently. The proposed marking scheme is not a decision of the team to harmonise marking.
